\documentclass[11pt,a4paper]{article}

\usepackage[margin=2.2cm]{geometry}
\usepackage{booktabs}
\usepackage{array}
\usepackage{tabularx}
\usepackage{amsmath}
\usepackage{amssymb}
\usepackage[hidelinks]{hyperref}
\usepackage{enumitem}

\newcolumntype{L}{>{\raggedright\arraybackslash}p{0.16\textwidth}}
\newcolumntype{C}{>{\raggedright\arraybackslash}p{0.30\textwidth}}
\newcolumntype{R}{>{\raggedright\arraybackslash}p{0.44\textwidth}}

\newcommand{\code}[1]{\texttt{#1}}

\title{PAID versus PhysVAE:\\
A comparison of the physics-aligned interpretable diffusion model\\ and the (now fixed) PhysVAE reimplementation\\
\large With a ranked diagnosis of PAID-specific issues ahead of the PAID epic}
\author{Generated by Mistral Vibe}
\date{\today}

\begin{document}
\maketitle

\section*{Scope and sources}

This report compares the two physics-integrated models that live side by side in the
private \code{interpretable-diffusion} repository at \code{main} commit
\code{d07ce2e} (``Aligning the reimplementation with the original (\#27)'', which
merged the PhysVAE diagnosis-and-fix epic):

\begin{itemize}[nosep]
  \item \textbf{PhysVAE} --- \code{src/models/phys\_vae.py},
        \code{src/models/train\_phys\_vae.py}, config
        \code{configs/diffusion\_vae/26MAR30.yaml}. This is the reference point:
        the advdif reimplementation was diagnosed in
        \code{reports/advdif\_reimplementation\_comparison.tex} (items D1--D8) and
        subsequently fixed; \code{main} now carries the paper recipe (x\_grid 20,
        integral\_level 2, $\mathit{x\_log\_sigma} = -13.8$, $\gamma = 10^6$,
        weight decay, grad clipping, best-validation checkpointing).
  \item \textbf{PAID} (physics-aligned interpretable diffusion) ---
        \code{src/models/paid.py} (\code{PAID} + \code{PhysUNet2D} +
        \code{SinusoidalEmbedding}), \code{src/models/train\_paid.py}
        (\code{paid\_loss}), \code{src/models/ddim.py} (\code{compute\_betas}),
        \code{src/visualisers/ddim.py} (\code{generate\_images}),
        \code{src/eval.py} (\code{evaluate\_paid}), and the committed configs
        \code{configs/diffusion\_paid/*.yaml}, \code{configs/pendulum\_paid/*.yaml}.
\end{itemize}

Per the epic handover (\code{PAID\_EPIC\_CONTEXT.md}), PAID is this repository's own
contribution --- a diffusion-based successor to PhysVAE --- so its reference material
is the project's own papers, not the original \code{phys-vae} repository; the
PhysVAE code here serves as the in-repo baseline whose conventions and fixes PAID
should inherit. Section~\ref{sec:diagnosis} ranks the PAID-specific gaps and bugs
found by static analysis. No training or evaluation runs were executed for this
report; every item states what would confirm it empirically.

\section{Setup at a glance}

\begin{table}[h]
\centering
\small
\begin{tabularx}{\textwidth}{LXX}
\toprule
 & \textbf{PhysVAE (fixed)} & \textbf{PAID} \\
\midrule
Paradigm & Variational autoencoder: encode $\to$ sample latents $\to$ ODE decode; ELBO + three regularizers & Conditional denoising diffusion: predict the clean signal $x_0$ from $(x_\tau, \tau)$ at each diffusion step; MSE + three regularizers \\
Latents & Stochastic: $(\mu, \log\sigma)$ per auxiliary and physical latent, prior + KL & Deterministic point estimates; no distributions, no priors, no KL \\
Known physics & $a\,\partial^2 y/\partial s^2$ in-ODE, advection via a \code{pre\_func} network & Same structure inside \code{PhysUNet2D.\_generate}: \code{physics\_model} + optional \code{pre\_funcs} \\
Coefficient estimate & Posterior mean $\mu$ of the physics latent (\code{phys\_vae.py:566--577}) & Output of a \code{sigmoid} + affine \code{phys\_func} head on the unmixer signal; at evaluation, read off the DDIM sampling trajectory \\
Data layout & $(B, T, L)$ flattened to $(B, L\cdot T)$ inside the model & $(B, T, X, 1)$ flattened to $(B, X\cdot T)$ inside the model \\
Conditioning & None (full reconstruction) & First \code{condition} frames kept clean; only the rest is diffused (\code{train\_paid.py:145--150}) \\
\bottomrule
\end{tabularx}
\end{table}

\section{Architecture}

\begin{table}[h]
\centering
\small
\begin{tabularx}{\textwidth}{LXX}
\toprule
\textbf{Component} & \textbf{PhysVAE} & \textbf{PAID} \\
\midrule
Encoder & Dedicated \code{aux\_funcs} $\to (2\,Z)$ split into $(\mu, \log\sigma)$; unmixer = \code{Linear(L,L)} + residual MLP (+ optional skip); \code{phys\_funcs} $\to (2\,Z)$ split & Monolithic denoiser \code{PhysUNet2D}: \code{aux\_funcs} $\to Z$ (deterministic); unmixer = single MLP $[Z + X\cdot T + E] \to X\cdot T$; \code{phys\_funcs} $\to Z$ via \code{sigmoid} + affine \\
Coefficient head & Linear $(\mu, \log\sigma)$ split; sample clipped to the model's range; evaluation uses the posterior mean & \code{sigmoid}$\times((5\,\mathrm{max}) - (\mathrm{min}/5)) + \mathrm{min}/5$: range $[\mathrm{min}/5,\,5\,\mathrm{max}]$, e.g.\ $[0.002, 0.5]$ for $[0.01, 0.1]$ (\code{paid.py:139--148}) \\
Time information & $\sin t, \cos t$ features inside the in-ODE \code{pre\_funcs} & Diffusion-step embedding: \code{SinusoidalEmbedding}($\code{embed\_dim}$) + 2-layer MLP, concatenated into every head \\
Embedding width & --- & \code{embed\_dim} = \code{data\_dim} = $\code{x\_grid}\times\code{t\_grid}$ (600 for advdif), so \code{time\_embed} alone is Linear(600, 2400)$+$Linear(2400, 600) and every \code{aux\_func} input is $600 + 600$ wide (\code{paid.py:96--116}) \\
Physics signal & Encoder unmixer output, supervised by \code{physics\_only\_loss} & Denoiser unmixer output, supervised by \code{phys\_loss} with per-frame $\alpha_\tau$ weighting in the noised region \\
ODE solve & One solve per decode; \code{integral\_level} substeps per output frame & Same solver machinery, shared \code{gen\_selections} combination enumeration; solved inside \code{full\_forward} at \emph{every} diffusion step during training and at every DDIM sampling step during evaluation \\
Post-ODE aux & \code{post\_funcs} applied per combination & Present in code, but \textbf{silently never applied} (diagnosis P1) \\
Parameter cost & MLPs sized by config & First-layer width tied to $\code{data\_dim} + \code{embed\_dim}$; burgers-shaped input ($\code{x\_grid}\times\code{t\_grid}$ in the tens of thousands) makes \code{time\_embed} a multi-billion-parameter layer --- known infeasibility, out of scope \\
\bottomrule
\end{tabularx}
\end{table}

\section{Losses}

\begin{table}[h]
\centering
\small
\begin{tabularx}{\textwidth}{LXX}
\toprule
\textbf{Term} & \textbf{PhysVAE (paper recipe)} & \textbf{PAID} \\
\midrule
Reconstruction & $\sum_{T,L}(x - \hat{x}_{PAB})^2$, mean over batch; balanced by $e^{x\_\mathrm{log\_ \sigma}} \approx 10^{-6}$ & $\sum_{T,X}(x - \hat{x}_{\text{last}})^2$, mean over batch --- full-sum scale, and it \emph{includes the clean conditioned frames} (no KL to balance) \\
KL / latents & KL to priors, weight $e^{-13.8}$ & None (deterministic) --- correctly absent from a diffusion objective \\
Unmixing (\code{phys\_only\_loss} / \code{phys\_loss}) & $\beta = 10^{-2}$, encoder signal vs \code{stop\_gradient} of the physics-only decode & \code{phys\_reg} $= 10^{-2}$, denoiser signal vs \code{stop\_gradient} of $x_{\text{physics-only}}$, with per-frame $\alpha_\tau$ weighting in the noised region (\code{train\_paid.py:124--139}) --- analogous and reasonable \\
Discrepancy & $\alpha = 10^{-1}$, per-sample Hamming-1 combination pairs (D5 fix shared via \code{gen\_mask}) & \code{disc\_reg} $= 10^{-1}$, same fixed per-sample implementation over $(m, B, T, X, 1)$ --- inherits the D5 fix \\
Data augmentation & $\gamma = 10^{6}$: random ground-truth coefficient $\to$ \code{generate\_physics} $\to$ regress the posterior \emph{mean} $\mu$ only (D4 fix) & \code{aug\_reg} $= 10^{4}$: same construction, and the scalar sigmoid head needs no $\mu/\log\sigma$ split --- D4 is structurally avoided, but the weight is $100\times$ smaller than the recipe (P4) \\
Timestep sampling & --- & Antithetic pairs: $\tau$ and $\code{num\_timesteps} - 1 - \tau$ (\code{train\_paid.py:134--137}) \\
Weight names & $\alpha$ / $\beta$ / $\gamma$ (paper naming) & \code{disc\_reg} / \code{phys\_reg} / \code{aug\_reg} --- semantics align with $\alpha/\beta/\gamma$ respectively, but names and the $\gamma$ value diverge \\
\bottomrule
\end{tabularx}
\end{table}

\section{Training and evaluation protocol}

\begin{table}[h]
\centering
\small
\begin{tabularx}{\textwidth}{LXX}
\toprule
\textbf{Aspect} & \textbf{PhysVAE} & \textbf{PAID} \\
\midrule
Trainer & Generic \code{train.py}: weight decay $10^{-6}$, value grad-clip 5.0, no-grad validation, best-validation checkpointing keyed on aux index 1 & Identical for free; \code{paid\_loss} aux order is $[\text{full}, \text{recon}, \text{phys}, \text{disc}, \text{aug}]$, so index 1 is reconstruction --- correct \\
Diffusion constants & --- & 1000 linear steps, $\beta \in [10^{-4}, 0.02]$ (\code{compute\_betas}), conditioning \code{condition} frames \\
Coefficient metric & \code{coeff\_mae} of the posterior mean vs true coefficient & \code{coeff\_mae} of the coefficient channel produced on the DDIM trajectory (\code{estimation\_type="mean"}, jump $= 10$), read at frame 0 (\code{eval.py:39--50}) \\
Forecast metric & MAE of the full decode & MAE of the DDIM-generated series, conditioned frames included in the average \\
Sampling cost & One ODE solve per datapoint & \code{full\_forward} (all $m$ combination ODE solves, unrolled through \code{diffeqsolve}) at each of the $\lceil 999/10 \rceil$ DDIM steps per batch \\
\bottomrule
\end{tabularx}
\end{table}

\section{Inheritance of the D1--D8 fixes}

The structural fixes landed during the PhysVAE epic are shared code, and PAID
inherits most of them by construction:

\begin{itemize}[nosep]
  \item \textbf{D1 (time-grid mismatch):} fixed. \code{PAID.\_\_init\_\_} derives
        $\code{dt} = \code{time}/(\code{t\_grid}-1)/\code{integral\_level}$ and
        $\code{t\_span}$ from \code{dataset\_config}
        (\code{paid.py:438--447}), exactly the PhysVAE convention.
  \item \textbf{D5 (discrepancy bug):} fixed. Both models share
        \code{gen\_selections}/\code{gen\_mask} and per-sample Hamming-1 pairs.
  \item \textbf{D4 (augmentation broadcasting):} structurally avoided --- the
        \code{phys\_func} head outputs a single scalar per range, so there is no
        $(\mu, \log\sigma)$ to mis-slice.
  \item \textbf{D6 (sampled latent at evaluation):} moot --- there is no posterior
        to sample; but see P6 on what the evaluation reads instead.
  \item \textbf{D7 (protocol):} inherited generically by \code{train.py}
        (weight decay, clipping, validation isolation, best-val checkpointing);
        the aux-index check in Section~above confirms it selects reconstruction.
  \item \textbf{D8 (reconstruction scale):} aligned --- full $\sum$ over
        $(T, X)$, mean over batch, matching the fixed PhysVAE convention.
\end{itemize}

\section{Ranked diagnosis of PAID-specific issues}
\label{sec:diagnosis}

\subsection*{P1 --- Post-ODE auxiliary functions are silently never applied}

In \code{PhysUNet2D.\_generate}, the post-function loop zips over
\code{aux\_latents[len(self.pre\_funcs) : len(self.post\_funcs)]}
(\code{paid.py:433--437}). With one pre-function and one post-function this slice is
\code{latents[1:1]} --- empty --- so \code{zip} terminates immediately and
\code{post\_funcs} never runs. The fixed PhysVAE slices
\code{latents[len(pre\_funcs) : len(post\_funcs) + 1]}
(\code{phys\_vae.py:363}), i.e.\ index 1. Consequence: any configured
\code{post\_funcs} is dead weight (parameters still created, gradients zero), and
the ``all auxiliary'' combination \code{x\_hat[-1]} used as the reconstruction
target is actually the physics+pre-only output. The committed advdif configs use
empty \code{post\_funcs}, so nothing currently observes it, but it makes the
generalization illusory. \textbf{Fix:} use the PhysVAE slice
(\code{len(post\_funcs) + 1}). Related cosmetics in the same function: the default
\code{post\_selections} uses \code{len(self.pre\_funcs)}
(\code{paid.py:248--252}) and one assert error message references
\code{pre\_funcs} for the post check.

\subsection*{P2 --- Committed PAID configs do not construct the current model}

\code{PAID.\_\_init\_\_} accepts \code{pde: \{integral\_level, solver,
phys\_param\_ranges\}} plus flat \code{aux\_funcs}/\code{unmixer}/
\code{phys\_funcs}/\code{pre\_funcs}/\code{post\_funcs} entries with \code{dims}
lists (the \code{make\_mlp} interface), and takes \code{condition}, \code{pad},
diffusion constants, and the three regularizer weights at the top level.
The committed configs violate this in several ways:

\begin{itemize}[nosep]
  \item \code{26MAY06.yaml} carries a dead \code{unet:} block (\code{PAID} has no
        \code{unet} parameter $\to$ \code{TypeError} on load) and
        \code{features}/\code{in\_features}-style \code{aux\_funcs}/
        \code{phys\_funcs} instead of \code{dims}.
  \item \code{mlp.yaml}/\code{pipeline\_test.yaml}/\code{condition\_base.yaml}
        carry a legacy \code{pde:} block with \code{system}, \code{timestep},
        \code{spatial\_length}, \code{num\_grid\_points}, \code{time\_dim} ---
        all now derived from \code{dataset\_config} and rejected by the
        constructor.
  \item \code{condition\_base.yaml} contains the literal placeholder
        \code{[CONDITION]}.
  \item \code{integral\_level: 1.0} is a float in all committed configs while the
        PhysVAE paper recipe uses 2 (integer).
  \item \code{trainer\_config} in the committed configs lacks the inherited
        \code{weight\_decay}/\code{grad\_clip} keys (they default off via
        \code{.get}, so the D7 fixes are silently inactive for PAID runs unless
        the config is rewritten).
\end{itemize}

This was left deliberately stale after the user reverted an earlier alignment
commit, so rewriting the configs (plus an uncommitted
\code{configs/diffusion\_paid/test\_diffusion.yaml}-style smoke config mirroring
the PhysVAE one) is the expected first plan item, not a unilateral fix.

\subsection*{P3 --- Coefficient head range $[\mathrm{min}/5,\,5\,\mathrm{max}]$}

The \code{phys\_func} head maps sigmoid output to
$x\,((5\,\mathrm{max}) - \mathrm{min}/5) + \mathrm{min}/5$
(\code{paid.py:139--148}). For the advdif range $[0.01, 0.1]$ this allows
$[0.002, 0.5]$ --- a range $25\times$ wider than the data range and far wider than
the PhysVAE model range $[\mathrm{min}/2,\,2\,\mathrm{max}] = [0.005, 0.2]$ that
the fixed PhysVAE uses for both its prior and clipping. Consequences: the head can
spend its sigmoid budget far outside the physically meaningful region, and the
augmentation regularizer (which draws ground truth from the same
\code{phys\_param\_ranges}) can supervise on coefficients the data never
exhibits. If this is not a deliberate PAID-paper choice, it is likely a
transcription slip of the half-min-to-double-max convention. \textbf{Fix:} confirm
against the PAID paper; if not intentional, use the PhysVAE range.

\subsection*{P4 --- Regularizer weights and naming}

\code{disc\_reg}/\code{phys\_reg}/\code{aug\_reg} map onto the paper's
$\alpha/\beta/\gamma$ semantics ($1/10$, $1/100$, augmentation), and the committed
values $10^{-1}/10^{-2}/10^{4}$ put the augmentation term $100\times$ below the
PhysVAE paper recipe's $\gamma = 10^{6}$. The scales of the two reconstruction
terms differ (PAID balances against a full-strength MSE, PhysVAE against an
ELBO), so the transfer is not automatic --- but the augmentation regularizer is
the only direct supervision of the coefficient, and the paper's own ablation says
coefficient inference collapses without it. This needs an explicit audit once the
weight semantics are confirmed against the PAID paper, plus a rename to
$\alpha$/$\beta$/$\gamma$ if the user opts for paper naming.

\subsection*{P5 --- Reconstruction trains on the conditioned frames}

\code{reconstruction\_loss} compares \code{data} to \code{x\_hat[-1]} over all
frames, including the \code{condition} frames that were fed to the denoiser clean
(\code{train\_paid.py:120--124}). The denoiser can satisfy those frames by
copying, and the ODE solve is anchored on \code{x[:, 0]} anyway, so the extra
terms mostly reward identity behaviour; PhysVAE reconstructs everything by
design, so this is a PAID-specific question of intent (the comment says
``improve smoothness''). Worth an explicit decision rather than a silent
default.

\subsection*{P6 --- Coefficient estimate rides on the DDIM trajectory}

\code{evaluate\_paid} runs 100-step DDIM sampling with
\code{estimation\_type="mean"} and reads the coefficient from the final step's
\code{phys\_params} channel at frame 0 (\code{eval.py:39--50},
\code{visualisers/ddim.py:143--159}). The value is therefore a byproduct of a
stochastic sampling trajectory evaluated on a nearly-denoised input, not a
dedicated inference of the coefficient from the observed series --- closer to
PhysVAE's old sampled-latent evaluation (D6) than to its fixed posterior-mean
protocol. If the intended PAID protocol is ``estimate $a$ from the observed
data'', calling the denoiser once on the (noised) observation at a fixed $\tau$
would be cheaper and less noisy. Unverified against the PAID paper --- this is
flagged per the epic context's D6-style open question.

\subsection*{P7 --- Compute cost of the wide embedding and per-step ODE solves}

\code{embed\_dim} = \code{data\_dim} makes every head input
$2\,\code{data\_dim}$-wide and \code{time\_embed} a
$600 \times 2400$-parameter MLP for advdif; combined with an \code{m}-combination
ODE solve inside every diffusion training step and every DDIM sampling step,
PAID's step cost is far above PhysVAE's. Not a correctness bug, but it caps grid
resolution (and is the root of the known burgers infeasibility, which remains out
of scope). A decoupled, smaller \code{embed\_dim} would remove the infeasibility
class if the PAID paper permits it.

\section{Suggested fix order}

\begin{enumerate}
  \item Rewrite the PAID configs to the current constructor (P2), including a
        committed full-run config mirroring \code{26MAR30.yaml}'s dataset recipe
        and the inherited \code{weight\_decay}/\code{grad\_clip} keys, plus an
        uncommitted smoke config mirroring
        \code{configs/diffusion\_vae/test\_diffusion.yaml}.
  \item Fix the post-function latent slice (P1) --- a two-character change with a
        \code{post\_funcs} smoke test proving the combination count and output
        actually change.
  \item Confirm the coefficient-head range against the PAID paper and align with
        the PhysVAE range if unintended (P3).
  \item Audit the three regularizer weights against the PAID paper; decide the
        $\alpha/\beta/\gamma$ rename and the augmentation weight (P4).
  \item Decide the reconstruction target over conditioned frames (P5) and the
        coefficient-evaluation protocol (P6) with the user --- both are
        paper-intent questions, not mechanical fixes.
  \item Optionally decouple \code{embed\_dim} from \code{data\_dim} (P7).
\end{enumerate}

Items 1--2 are mechanical and unblock everything else; items 3--6 are user
decisions per the epic's working agreements (plain questions, no suggested
options, PLAN.md blocks with one commit each). The ranking is from static
analysis; P1 is trivially confirmable by counting applied combinations in a
one-batch smoke run, and P2 by simply loading each committed config.

\end{document}
